\section{A support criterion at a specified index}
\label{ext:direct-gaussian-support}
The support calculations throughout the finite range use one elementary
criterion. Its inputs are a lower bound for the attained score, upper
bounds for moments of the same marginal, and a predecessor allowance.
The Gaussian calculation keeps the actual integer $n$ throughout.

\begin{lemma}[Direct Gaussian support criterion]
\label{ext:direct-gaussian-support-criterion}
Let $n\ge2$, and let $(P,H)$ be an attained positive closed-event
ratio maximizer with score $C$. Write
$\rho=\E|X|^3$, $\beta=\E[X|X|]$ and
$A=\sqrt{n-1}\Delta_{n-1}(P)/\rho$.
Suppose known numbers $c,r,b_-,b_+,S,e,q$ satisfy
\[
 0<c\le C,\quad 1\le\rho\le r,\quad
 -1<b_-\le\beta\le b_+<1,\quad q\ge0,
\]
\[
 \sqrt{n/(n-1)}\le S,\qquad
 n(\sqrt{n/(n-1)}-1)-1/2\le e,\qquad
 n(A-C)_+\le q.
\]
Fix $R\ge1$ and a support sign $\eta\in\{-1,1\}$, and put
$B_+=b_+$, $B_-=-b_-$. Define
\[
 Q_\eta(v)=-v^3+\tfrac32rv^2+3B_\eta v+re,
 \qquad
 L_\eta(v)=\tfrac25(v^2/2+1/2+e)+cQ_\eta'(v).
\]
Suppose that
\begin{equation}\label{ext:direct-support-ray-gates}
 Q_\eta(R),Q_\eta'(R),Q_\eta''(R),
 L_\eta(R),L_\eta'(R)<0,
 \qquad 1/5-3c<0.
\end{equation}
Define the actual-index Gaussian expression
\[
 G_n(H,t)=n\sqrt n\{\Phi((H-t)/\sqrt{n-1})-\Phi(H/\sqrt n)\}
 +nt\phi(H/\sqrt n)+\tfrac H2(t^2-1)\phi(H/\sqrt n),
\]
and suppose $U_\eta$ satisfies
\[
 \sup_{H\in\R}\eta G_n(H,R)\le U_\eta,
 \qquad U_\eta+cQ_\eta(R)+rSq<0.
\]
Then $\supp P$ contains no point in the ray
$\{\eta v:v\ge R\}$. The two signs can be checked separately,
with different signed-moment bounds.
\end{lemma}

\begin{proof}
Contact gives
\[
 0\le G_n(H,t)+
 C\{-|t|^3+\tfrac32\rho t^2+3\beta t+\rho e_n\}
 +\rho s_n n(A-C),
\]
where $s_n=\sqrt{n/(n-1)}$ and $e_n=n(s_n-1)-1/2$.
The last term is at most $rSq$. At $t=\eta v$ with $v\ge R$,
the bracket is at most $Q_\eta(v)$.
Since $Q_\eta'''=-6$, the first three gates make this polynomial
strictly negative on the full ray. Thus multiplication by $C\ge c$
gives an upper bound $cQ_\eta(v)$.

The derivative is
$n\phi(z)+Ht\phi(z)-ns_n\phi(s_ny)$, where
$z=H/\sqrt n$ and $y=(H-t)/\sqrt n$.
Taylor's theorem and $\|\phi''\|_\infty=\phi(0)$ bound
$|n\phi(y)-n\phi(z)-Ht\phi(z)|$ by $\phi(0)t^2/2$.
The identity $\partial_s[s\phi(sy)]=-\phi''(sy)$ bounds the
remaining scaling error by $n\phi(0)(s_n-1)$.
This proves, for every $n\ge2$,
$|\partial_tG_n(H,t)|\le\phi(0)(t^2/2+1/2+e)$.
Therefore the derivative in $v$ of
$G_n(H,\eta v)+cQ_\eta(v)+rSq$ is at most $L_\eta(v)$.
Its second derivative as a polynomial bound is
$L_\eta''=2/5-6c<0$; the last two endpoint gates imply
$L_\eta(v)<0$ on $[R,\infty)$.
Finally $G_n(H,-R)=-G_n(-H,R)$, so the assumed $U_\eta$
bounds $G_n(H,\eta R)$ uniformly in $H$.
The strict endpoint inequality and monotonicity contradict contact
on every point of the specified support ray.
\end{proof}

\paragraph{How the Gaussian-line certificate supplies $U_+$ and $U_-$.}
Let $z=H/\sqrt n$ and $a=(H-R)/\sqrt{n-1}$. Direct
differentiation gives
\[
 \partial_HG_n(H,R)=ns_n\phi(a)+
 \left[-n-RH+\tfrac12(R^2-1)(1-H^2/n)\right]\phi(z).
\]
Choose an integer $H_0$ such that
$H_0\ge5\sqrt n$ and $H_0-R\ge5\sqrt{n-1}$.
An outward enclosure of the square roots makes this choice exact.
On the whole real line,
\[
 |\partial_H^2G_n(H,R)|\le M_{n,R}
 :=\frac{\sqrt n}{4}\left(\frac n{n-1}+1\right)
   +\frac25R+\frac{2|R^2-1|}{5\sqrt n}.
\]
On a closed threshold cell with midpoint $h$ and half-width $w$,
let $[g_-,g_+]$ and $[d_-,d_+]$ enclose $G_n(h,R)$ and
$\partial_HG_n(h,R)$. Taylor's theorem then gives the two upper
bounds
\[
 g_++\max(|d_-|,|d_+|)w+M_{n,R}w^2/2,
 \qquad
 -g_-+\max(|d_-|,|d_+|)w+M_{n,R}w^2/2.
\]
They bound $G_n(H,R)$ and $-G_n(H,R)$ respectively throughout
the cell. The certificate includes every consecutive cell from
$-H_0$ to $H_0$, including their endpoints.

On the two exterior rays, both distribution-function arguments lie
in the same normal tail, so
\[
 |G_n(H,R)|\le n\sqrt n\,\tau+
 \left(nR+\tfrac12H_0|R^2-1|\right)\phi(H_0/\sqrt n),
 \quad |H|\ge H_0,
\]
where $\tau=3\cdot10^{-7}>\Phi(-5)$ is sufficient.
Take the largest positive-cell bound and the exterior bound for
$U_+$, and the largest negative-cell bound and the exterior bound
for $U_-$. This proves a statement about the entire threshold line,
not just about its grid points.

\paragraph{Exact arithmetic and the native input interface.}
The implementation \texttt{direct\_gaussian\_support.py} uses
integer-bracketed square roots and outward rational interval operations.
Its density enclosure first divides the nonnegative exponential argument
by a power of two until it is at most one, uses the alternating Taylor
polynomial through degree23 with the degree24 upper remainder, and
repeatedly squares with outward rounding. For a nonnegative normal
argument at most five, its distribution-function enclosure integrates
the first80 terms of the exponential series and adds the next-term
error; floor and ceiling arithmetic encloses every term. The normal
argument is rounded outward to a grid of width $10^{-12}$, and
monotonicity of $\Phi$ supplies the interval endpoints. Outside this
range it uses the preceding Mills bound. The density constant is
enclosed by Machin's arctangent identity and integer square roots.
These rules justify the midpoint value and derivative intervals used
above.

The native function \texttt{target(n,K,R,q,...)} uses
$c=819/2000+K/n$. The hypothesis $C\ge\CE+K/n$ implies
$C>c$. Its default third-moment cap comes from the two published
affine Berry--Esseen estimates already used in the introduction:
\[
 r=\min\left\{
 \frac{0.3031\cdot0.646}{c-0.3031},
 \frac{0.3322\cdot0.429}{c-0.3322}\right\}.
\]
Both denominators are positive on its admitted domain. The signed-moment
inequality gives the default symmetric interval for $\beta$.
An explicitly supplied tighter $\rho$ or signed-$\beta$ interval must
come from a previously established property of the same marginal.
The program checks all the strict polynomial gates and the two full
Gaussian-line bounds. Its return value proves the support conclusion
of this lemma only after these analytic and same-marginal hypotheses
have been supplied. In particular, it does not itself prove a score
bound or transfer a result from one index to another.
